\subsection{Ablation Study across different level of contamination}

Expanding on the ablation table in the main paper, in Tab.~\ref{tab:results_cifar10_ablation_all_zeta} we show ablation comparisons (augmentation only, gradient step modification only, or both) across all tested values of the mismatch in labeled-vs-unlabeled class content $\zeta$.

\begin{table}[!h]
	\setlength{\tabcolsep}{.08cm}
    \begin{tabular}{c c}
		{\large Mismatch $\zeta=0\%$}
		&
		{\large Mismatch $\zeta=25\%$}
		\\
		\begin{minipage}{.5\textwidth}
		\resizebox{\textwidth}{!}{%	
			\begin{tabular}{r | r | r | r | r}
    & ~~~~ & +A only & +G only & +A\&G (Fix-A-Step)
\\ 
    Pi-Model & 78.32 & 81.90 & 78.48 & \textbf{82.83} 
\\ 
    Mean-Teacher & 79.57 & \textbf{84.18} & 80.60 & 83.02
\\ 
    VAT & 79.15 & \textbf{83.88} & 79.47 & 83.10
\\
    Pseudo-label & 77.43 & 79.03 & 78.30 & \textbf{79.98}
\\
    FixMatch & 86.40 & 86.35 & 86.17 & \textbf{88.00}
\end{tabular}	
		}
		\end{minipage}
		&
		\begin{minipage}{.5\textwidth}
		\resizebox{\textwidth}{!}{%	
			\begin{tabular}{r | r | r | r | r}
             & ~~~~ & +A only & +G only & +A\&G (Fix-A-Step)
\\ 
    Pi-Model & 77.45 & 79.12 & 78.00 & \textbf{79.80} 
\\ 
    Mean-Teacher & 77.70 & 81.35 & 78.28 & \textbf{87.77}
\\ 
    VAT & 76.65 & 82.27 & 78.35 & \textbf{82.63}
\\
    Pseudo-label & 76.58 & 79.28 & 77.25 & \textbf{79.60}
\\
    FixMatch & 83.20 & 84.44 & 83.85 & \textbf{86.63}
\end{tabular}	
		}
		\end{minipage}
		\\ %%%%%%%%%%%%%%%%%%%%%%%%%%
		\\ %%%%%%%%%%%%%%%%%%%%%%%%%%
		{\large Mismatch $\zeta=50\%$}
		&
		{\large Mismatch $\zeta=75\%$}
		\\
		\begin{minipage}{.5\textwidth}
		\resizebox{\textwidth}{!}{%	
			\begin{tabular}{r | r | r | r | r}
             & ~~~~ & +A only & +G only & +A\&G (Fix-A-Step)
\\ 
    Pi-Model & 76.03 & 78.70 & 76.57 & \textbf{78.97} 
\\ 
    Mean-Teacher & 76.35 & \textbf{82.22} & 78.18 & 80.15
\\ 
    VAT & 75.90 & 79.43 & 77.37 & \textbf{79.56}
\\
    Pseudo-label & 75.42 & 78.33 & 75.65 & \textbf{78.77}
\\
    FixMatch & 81.60 & 83.28 & 81.80 & \textbf{85.45}
\end{tabular}	
		}
		\end{minipage}
		&
		\begin{minipage}{.5\textwidth}
		\resizebox{\textwidth}{!}{%	
			\begin{tabular}{r | r | r | r | r}
             & ~~~~ & +A only & +G only & +A\&G (Fix-A-Step)
\\ 
    Pi-Model & 75.00 & 77.82 & 74.77 & \textbf{79.35} 
\\ 
    Mean-Teacher & 74.33 & 79.63 & 74.52 & \textbf{81.08}
\\ 
    VAT & 74.87 & 80.82 & 75.20 & \textbf{81.20}
\\
    Pseudo-label & 76.85 & 78.38 & 77.15 & \textbf{78.83}
\\
    FixMatch & 81.05 & 83.03 & 81.33 & \textbf{84.75}
\end{tabular}	
		}
		\end{minipage}
	\end{tabular}
	\caption{\textbf{Ablation analysis on CIFAR-10 6 animal task}, examining how accuracy changes for each SSL method if we only use our augmentation (+A), only use our gradient step modification (+G), and use the combination (+A\&G) which constitutes our Fix-A-Step. 
	Each panel shows results for a fixed value of the mismatch percentage $\zeta$ describing the overlap in classes between labeled and unlabeled set.
    For each method, we \textbf{bold} the best result.
    Setting: 400 examples/class.
    }%endcaption
\label{tab:results_cifar10_ablation_all_zeta}
\end{table}


\subsection{Robustness of results across multiple train/test splits}

In the main paper, we report results on CIFAR-10 6 animal task across many baselines methods.
For each baseline, we use only one train/test split due to the huge computation required to compare all baselines.
In Fig.~\ref{fig:cifar10_uncertainty}, for a subset of methods we assess the \emph{robustness} of the conclusions of that experiment across multiple separate training/test splits.

We train the labeled-set-only baseline, FixMatch with and without Fix-A-Step, and Mean-Teacher with and without Fix-A-Step for 5 random splits of the data, across two levels of mismatch ($\zeta = 50\%$ and $\zeta = 100\%$).
Results are in Fig.~\ref{fig:cifar10_uncertainty}.
Broadly, we suggest that our conclusion that Fix-A-Step delivers successful accuracy gains holds even across 5 splits: both FixMatch and MeanTeacher show notable gains across both levels of mismatch $\zeta$.
\textbf{Notably, MeanTeacher plus Fix-A-Step appears quite competitive with off-the-shelf FixMatch}, and FixMatch plus Fix-A-Step is the best of all.


\begin{figure*}[!h]
\centering
\begin{tabular}{c c}
\includegraphics[width=.35\textwidth]{figures/uncertainty_cifar10_400perclass/OliverContaminationLevel50_StdErrorBar_ChangeFont.pdf}
&
\includegraphics[width=.35\textwidth]{figures/uncertainty_cifar10_400perclass/OliverContaminationLevel100_StdErrorBar_ChangeFont.pdf}
\end{tabular}
\caption{
\textbf{CIFAR10, 400 examples/class}. Bar height indicates the average across 5 random splits of the data. Error bars show the standard deviation across 5 splits.
}%endcaption	
\label{fig:cifar10_uncertainty}
\end{figure*}

\subsection{Sensitivity to hyperparameters}
Since Deep SSL methods could be sensitive to hyper-parameters, we conduct sensitivity analysis to see how Fix-A-Step behave under different choice of sharpening temperature $\tau$ and Beta distribution shape $\alpha$. We analyzed the performance of Fix-A-Step using a Mean-Teacher base model across several reasonable choices of MixUp parameter $\alpha {\in} \{0.5, 0.75\}$ and sharpening temperature $\tau {\in} \{0.5, 0.95\}$ (totally 4 combinations). (See. Alg.~\ref{alg:FixAStep} for hyperparameter definitions).
Results in Fig.~\ref{fig:cifar10_sensitivity} shows that Fix-A-Step's performance gains over the off-the-shelf (or ``vanilla'') Mean Teacher base do not appear overly sensitive to these hyperparameters.

\begin{figure*}[!h]
\centering
\begin{tabular}{c c}
\includegraphics[width=.42\textwidth]{figures/sensitivity_cifar10_400perclass_MT/OliverContaminationLevel50_MT_ChangeFont.pdf}
&
\includegraphics[width=.42\textwidth]{figures/sensitivity_cifar10_400perclass_MT/OliverContaminationLevel100_MT_ChangeFont.pdf}
\end{tabular}
\caption{
\textbf{CIFAR10, 400 examples/class}. \emph{Vanilla MT}: Mean teacher SSL base model. \emph{Fix-A-Step MT1}: MT base Fix-A-Step with $\tau=0.5$ , $\alpha=0.5$, . \emph{Fix-A-Step MT2}: MT base Fix-A-Step with $\tau=0.95$, $\alpha=0.75$. \emph{Fix-A-Step MT3}: MT base Fix-A-Step with $\tau=0.5$, $\alpha=0.75$. \emph{Fix-A-Step MT4}: MT base Fix-A-Step with $\tau=0.95$, $\alpha=0.5$.  Results average across 5 random split of the data. Error bar showing standard deviation across the 5 split.
}%endcaption	
\label{fig:cifar10_sensitivity}
\end{figure*}

\vspace{2cm}
\subsection{Comparison of computation cost and performance}
\setlength{\tabcolsep}{.3cm}
\begin{table}[!h]
\centering
\begin{tabular}{c | r r | r r | r r }
\multicolumn{0}{c}{Methods} &\multicolumn{2}{c}{$\zeta=25\%$} &\multicolumn{2}{c}{$\zeta=50\%$} &\multicolumn{2}{c}{$\zeta=75\%$}\\

& Acc & Runtime & Acc & Runtime & Acc & Runtime 
 \\ \hline
MT+Fix-A-Step & 81.77 & 1476 & 80.15 & 1331 & 81.08 & 1333\\
VAT+Fix-A-Step & 82.63 & 1662 & 79.56 & 1691 & 81.20 & 1681\\
OpenMatch & 85.45 & 2728 & 79.56 & 2803 & 79.62 & 3270\\
\end{tabular}
    \caption{
    \textbf{Comparison of runtime and test accuracy on CIFAR-10 6 animal task}. Setting: 400 example/class. Mismatch percentage $\zeta$ describes the overlap in classes between labeled and unlabeled set. Runtime (in minutes) is based on training same number of steps on a Nvidia A100 GPU. 
    }%endcaption
    \label{tab:cifar10_runtime_acc}
\end{table}
